% Bewertungspaket zum Gesamttest «Potenz- und Wurzelfunktionen» — Quelle des PDFs (python3 scripts/build-lp-pdf.py).
% Ganze Punkte je Teilschritt; Werte mit python3 nachgerechnet (03.10.2026, Folgefehler-Fälle eingeschlossen).
\documentclass[11pt]{article}
\usepackage{../lp-druck}
\renewcommand{\lpname}{Potenz- und Wurzelfunktionen}
\renewcommand{\lpbereich}{Schwerpunktfach}
\renewcommand{\lpteilgebiet}{3.2}
\renewcommand{\lpfuss}{Bewertungspaket}
\newcolumntype{P}{>{\centering\arraybackslash}p{10mm}}
\newenvironment{raster}{\par\smallskip\noindent\tabularx{\linewidth}{@{}>{\raggedright\arraybackslash}p{38mm}P X@{}}\toprule
  \small\textsc{Teilschritt} & \small\textsc{P} & \small\textsc{Musterlösung}\\\midrule}{\bottomrule\endtabularx\par}
\newcommand{\fehler}[1]{\par\smallskip{\small\textcolor{lprot}{\bfseries Typische Fehler:} #1}}
\newcommand{\E}{\,{\footnotesize\bfseries\color{lpgruen}(E)}}
\newcommand{\A}{\,{\footnotesize\bfseries\color{lpblau}(A)}}
\begin{document}
\lpkopf{Bewertungspaket zum Gesamttest}{}

\begin{lpachtung}{Erst lösen, dann öffnen}
Dieses Paket enthält die Musterlösung. Wer es vor dem Lösen liest, misst am Ende nur, wie gut das Abschreiben gelingt.
\end{lpachtung}

\section*{1 · So gehst du vor}
\begin{enumerate}[leftmargin=*,itemsep=2pt]
\item \textbf{Gesamttest lösen} — vollständig, mit Rechenweg, ohne dieses Paket und ohne Taschenrechner.
\item \textbf{Fotografieren oder scannen}, gut lesbar, eine Seite pro Bild. \textbf{Kein Name} auf den Bildern.
  Handyfotos tragen oft unsichtbar Standort, Zeit und Gerät mit (Metadaten). Lade darum lieber einen Scan oder ein
  Bildschirmfoto des Fotos hoch, oder entferne vorher die Standortangaben.
\item \textbf{Bewerten lassen.} Ob du dafür eine KI nutzt und welche, entscheidest du selbst und in eigener Verantwortung —
  je nachdem, was dir aufgrund deines Alters und deiner persönlichen Situation erlaubt ist (Altersgrenzen und
  Nutzungsbedingungen des Anbieters, allenfalls Einverständnis deiner Eltern). Lade nur deine Lösung und dieses PDF hoch,
  keine persönlichen Daten, und füge den Auftrag aus Abschnitt~2 ein. Ohne KI kannst du dich mit Abschnitt~3 selbst bewerten.
\item \textbf{Rückmeldung prüfen.} Stimmt, was die KI aus deiner Handschrift gelesen hat? Eine KI kann sich irren —
  im Zweifel gilt das Raster in Abschnitt~3.
\item \textbf{Gezielt wiederholen} nach Abschnitt~4.
\end{enumerate}

\needspace{14\baselineskip}
\section*{2 · Auftrag an die KI}
\begin{tcolorbox}[colback=lplinie!25,colframe=lplinie,boxrule=0.5pt,arc=1mm,fontupper=\small]
Du bist Korrektorin oder Korrektor für Mathematik an einer Schweizer Berufsmaturitätsschule. Im Anhang findest du (1)~das Bewertungspaket zum Gesamttest «Potenz- und Wurzelfunktionen» mit Musterlösung und Punkteraster und (2)~Fotos meiner handschriftlichen Lösung.\par\medskip
Geh so vor:\par
1. Schreib zu jeder Aufgabe G1 bis G8 zuerst ab, was du in meiner Lösung liest — Rechenweg und Ergebnis. Was du nicht sicher lesen kannst, markierst du mit [unsicher]. Nicht raten.\par
2. Vergib die Punkte genau nach dem Raster im Bewertungspaket (Abschnitt 3), ganze Punkte je Zeile. Ziehe einen Fehler nur einmal ab: Rechne ich mit einem falschen Zwischenergebnis richtig weiter, gibt es die Folgepunkte. Die Zeilen «Typische Fehler» sagen, welche Rasterzeile bei diesem Fehler 0 Punkte gibt — sie sind kein zusätzlicher Abzug.\par
3. Andere richtige Lösungswege zählen voll. Gleichwertige Schreibweisen zählen voll: Dezimalpunkt oder Dezimalkomma, Bruch oder Dezimalzahl ($\tfrac12 = 0.5$, $\sqrt[2]{x} = \sqrt{x}$), $y = \ldots$ oder $f(x) = \ldots$, umgestellte Terme ($x^{-2}$ statt $\tfrac{1}{x^{2}}$) — ausser wo das Raster eine bestimmte Form verlangt. Drei Stufen: Zeilen mit \textbf{(A)} verlangen nur Ablesen, diese Punkte gibt es auch ohne Rechenweg. Zeilen mit \textbf{(E)} sind Ergebnispunkte: Das Ergebnis muss stimmen \emph{und} aus einem erkennbaren Weg hervorgehen, ausser die Zeile sagt ausdrücklich etwas anderes. Alle übrigen Zeilen bewerten den Weg selbst.\par
4. Begründe jeden Abzug in einem Satz und nenne die Fehlerart (zum Beispiel «Vorzeichen von $u$ in der Klammer»). Schreib die Musterlösung nicht ab — zeig mir, wo mein Fehler liegt.\par
5. Gib zum Schluss eine Tabelle «Aufgabe | Punkte | Begründung», die Gesamtpunktzahl (von 26) und — nach der Selbsteinschätzung in Abschnitt 4 — welches Kapitel des Leitprogramms ich wiederholen soll. Keine Note.\par
6. Fehlt ein Foto oder ist eine Aufgabe unlesbar, sag es, statt sie zu bewerten.
\end{tcolorbox}

\needspace{16\baselineskip}
\section*{3 · Musterlösung und Punkteraster}
\textbf{Allgemein:} Ganze Punkte je Zeile. Ein Fehler wird nur einmal abgezogen (Folgefehler). Gleichwertige Wege und
Schreibweisen zählen voll (Bruch oder Dezimalzahl, Terme in anderer Reihenfolge). In der Spalte P:
\textbf{(A)}~= nur ablesen, zählt auch ohne Rechenweg · \textbf{(E)}~= Ergebnispunkt, Ergebnis richtig
\emph{und} Weg erkennbar (ausser die Zeile sagt ausdrücklich etwas anderes) · ohne Zeichen = die Zeile bewertet den Weg.
Der ganze Test ist ohne Taschenrechner gestellt; Lineal und Bleistift sind erlaubt.
«Typische Fehler» nennt, welche Zeile 0 Punkte gibt, und die Punkte, die bleiben. Steht dort ein Bereich
(«1–2 von 3\,P»), hängt der zweite Punkt an der Bedingung dahinter.
\textbf{Teilrichtige Zeilen:} Verlangt eine Zeile mehrere Angaben und stimmt nur ein Teil, gibt die Zeile 0 Punkte —
ausser die Zeile nennt ausdrücklich eine Aufteilung.

\begin{lpaufgabe}{G1}{$f(x) = 2x^{4}$}{3 P}
\begin{raster}
(a) $f(-2)$ & 1\E & $f(-2) = 2 \cdot (-2)^{4} = 2 \cdot 16 = 32$\\
(a) $f(0.5)$ & 1\E & $f(0.5) = 2 \cdot 0.0625 = 0.125$ (gleichwertig $\tfrac18$)\\
(b) Symmetrie mit Nachweis & 1 & $f(-x) = 2(-x)^{4} = 2x^{4} = f(x)$: gerade Funktion, Graph achsensymmetrisch zur $y$-Achse. Nachweis und Aussage nötig.\\
\end{raster}
\fehler{$f(-2) = -32$ (Klammer vergessen, also $-2 \cdot 2^4$): Zeile (a) erste 0 $\to$ 2 von 3\,P.
$f(-2) = -16$ (Exponent als Faktor gerechnet, $2 \cdot 4 \cdot (-2)$): dieselbe Zeile 0.
$f(0.5) = 1$ ($(2 \cdot 0.5)^4$ statt $2 \cdot 0.5^4$ — $a$ in die Basis gezogen): Zeile (a) zweite 0 $\to$ 2 von 3\,P.
Nur «achsensymmetrisch» ohne $f(-x)$: Zeile (b) 0 $\to$ 2 von 3\,P. «$n$ ist gerade, darum achsensymmetrisch» zählt als Nachweis voll.}
\end{lpaufgabe}

\begin{lpaufgabe}{G2}{Kurve $\to$ Gleichung}{3 P}
\begin{raster}
$a$ aus $(1 \mid -1)$ & 1\A & $f(1) = a \cdot 1^{n} = a$, abgelesen $-1$, also $a = -1$\\
$n$ aus $(2 \mid -8)$ & 1 & $-1 \cdot 2^{n} = -8 \Rightarrow 2^{n} = 8 \Rightarrow n = 3$\\
Gleichung und Begründung & 1\E & $f(x) = -x^{3}$, mit einem Satz zur Punktsymmetrie bzw. zum fallenden Verlauf als Begründung für den ungeraden Exponenten\\
\end{raster}
\fehler{$a = 1$ (Vorzeichen übersehen): $a$-Zeile 0; $n = 3$ aus $1 \cdot 2^n = -8$ ist nicht lösbar — wer trotzdem $n = 3$ schreibt, bekommt die $n$-Zeile $\to$ 2 von 3\,P.
$n = 4$ ($f(2) = -16$ statt $-8$ abgelesen): $n$-Zeile 0 $\to$ 2 von 3\,P.
$a$ und $n$ vertauscht ($f(x) = -3x^{1}$ o.\,ä.): beide Zeilen 0 $\to$ höchstens 1 von 3\,P.
Gleichung richtig, Begründung fehlt: letzte Zeile 0 $\to$ 2 von 3\,P.}
\end{lpaufgabe}

\begin{lpaufgabe}{G3}{$g(x) = -\tfrac{2}{x^{2}}$}{3 P}
\begin{raster}
(a) Definitionsmenge & 1\A & $D = \mathbb{R}\setminus\{0\}$ (gleichwertig: alle $x$ ausser $0$)\\
(a) Asymptoten & 1\A & $x = 0$ und $y = 0$ (gleichwertig: die beiden Achsen)\\
(b) Lage der Äste mit Begründung & 1 & Beide Äste liegen \emph{unten}: Der Exponent $2$ ist gerade, also ist $x^{2} > 0$ für alle $x \neq 0$; mit $a = -2 < 0$ wird $g(x)$ immer negativ. Aussage und Begründung nötig.\\
\end{raster}
\fehler{$D = \mathbb{R}$: erste Zeile 0 $\to$ 2 von 3\,P.
Nur eine Asymptote genannt: zweite Zeile 0 $\to$ 2 von 3\,P.
«diagonal» (Parität verwechselt): Zeile (b) 0 $\to$ 2 von 3\,P.
«beide oben» (Vorzeichen von $a$ übersehen): Zeile (b) 0 $\to$ 2 von 3\,P.
Richtige Lage ohne Begründung: Zeile (b) 0 $\to$ 2 von 3\,P. Ein Zahlenbeispiel wie $g(1) = g(-1) = -2$ zählt als Begründung voll.}
\end{lpaufgabe}

\begin{lpaufgabe}{G4}{$h(x) = (x-1)^{4} - 16$ · $g(x) = \tfrac{2}{x-3} + 1$}{4 P}
\begin{raster}
(a) Ansatz & 1 & $(x-1)^{4} = 16$\\
(a) Beide Nullstellen & 1\E & $x - 1 = \pm 2 \Rightarrow x_1 = -1$, $x_2 = 3$. Beide nötig.\\
(a) Begründung für die Anzahl & 1 & Der Exponent ist gerade: Beim Wurzelziehen gibt es die positive \emph{und} die negative Lösung, darum genau zwei.\\
(b) Definitionsmenge und Asymptoten & 1\A & $D = \mathbb{R}\setminus\{3\}$; $x = 3$ und $y = 1$. Alle drei Angaben nötig.\\
\end{raster}
\fehler{Nur $x = 3$ (das $\pm$ vergessen): Nullstellen-Zeile 0; eine Begründung, die dazu passt («nur eine»), ist falsch und zählt ebenfalls 0 $\to$ 1 von 3\,P.
$x_1 = 1$, $x_2 = -3$ (Vorzeichen beim Auflösen vertauscht): Nullstellen-Zeile 0 $\to$ 2 von 3\,P.
$\sqrt[4]{16} = 4$ gerechnet, also $x = 5$ und $x = -3$: Nullstellen-Zeile 0, Begründung kann voll zählen $\to$ 2 von 3\,P.
Nullstellen richtig, Begründung fehlt: dritte Zeile 0 $\to$ 3 von 4\,P.
Bei (b) nur eine Asymptote oder $D = \mathbb{R}$: Zeile (b) 0 $\to$ 3 von 4\,P.
Bei (b) $x = -3$ oder $y = 3$ (Vorzeichen in der Klammer bzw. Zahlen vertauscht): Zeile (b) 0 $\to$ 3 von 4\,P.}
\end{lpaufgabe}

\begin{lpaufgabe}{G5}{$f(x) = x^{3} - 2$ umkehren}{4 P}
\begin{raster}
(a) nach $x$ aufgelöst & 1 & $y = x^{3} - 2 \Rightarrow x^{3} = y + 2 \Rightarrow x = \sqrt[3]{y+2}$\\
(a) vertauscht & 1\E & $f^{-1}(x) = \sqrt[3]{x+2}$\\
(b) Bildpunkt & 1\A & $(6 \mid 2)$ — beim Spiegeln an $y = x$ tauschen die Koordinaten; zählt auch ohne Rechenweg\\
(c) Skizze & 1 & Beide Kurven und die gestrichelte Gerade $y = x$; die zweite Kurve ist erkennbar das Spiegelbild der ersten. Toleranz: an den Rändern höchstens eine halbe Kästchenbreite daneben.\\
\end{raster}
\fehler{$f^{-1}(x) = \sqrt[3]{x} + 2$ (die $2$ nicht unter die Wurzel genommen): Zeile (a) zweite 0; wenn (a) erste stimmt, bleibt deren Punkt $\to$ 3 von 4\,P.
$f^{-1}(x) = \sqrt[3]{x-2}$ (Vorzeichen): Zeile (a) zweite 0 $\to$ 3 von 4\,P.
$f^{-1}(x) = x^{-3} + 2$ (Kehrwert statt Umkehrfunktion): beide (a)-Zeilen 0 $\to$ höchstens 2 von 4\,P.
$(2 \mid 6)$ bei (b) (nicht getauscht): Zeile (b) 0 $\to$ 3 von 4\,P.
$(-6 \mid -2)$ (getauscht, aber zusätzlich beide Vorzeichen gewechselt — das wäre die Spiegelung an $y = -x$): Zeile (b) 0 $\to$ 3 von 4\,P.
Bei (c) fehlt die Gerade $y = x$: Skizze 0, wenn die zweite Kurve auch fehlt; ist die Spiegelkurve richtig gezeichnet, zählt die Zeile trotzdem.
Die Spiegelkurve passt zur eigenen, falschen Umkehrfunktion: Zeile (c) zählt voll (Folgefehler).}
\end{lpaufgabe}

\begin{lpaufgabe}{G6}{Warum einschränken?}{3 P}
\begin{raster}
$x^{4}$: Begründung & 1 & $4$ ist gerade, also $f(-x) = f(x)$: Zu einem $y$ gehören zwei $x$. Beim Spiegeln lägen über einem $x$ zwei Punkte — das wäre kein Funktionsgraph.\\
$x^{4}$: Zahlenbeispiel & 1 & z.\,B. $(-2)^{4} = 2^{4} = 16$: Zu $y = 16$ gehören $x = -2$ und $x = 2$.\\
$x^{5}$: Begründung mit Beispiel & 1 & $5$ ist ungerade, die Kurve steigt überall: Jeder Wert kommt genau einmal vor, z.\,B. $(-2)^{5} = -32 \neq 32 = 2^{5}$.\\
\end{raster}
\fehler{Nur «der Exponent ist gerade» ohne die Folge für die Umkehrung: erste Zeile 0 $\to$ höchstens 2 von 3\,P.
Zahlenbeispiel fehlt oder ist keins (etwa «weil die Parabel symmetrisch ist»): die betroffene Zeile 0.
Die Einschränkung mit der Definitionsmenge der \emph{Wurzel} begründet («$\sqrt{x}$ gibt es erst ab $0$»): Das ist eine Folge, keine Begründung — erste Zeile 0 $\to$ höchstens 2 von 3\,P.
Beide Fälle vertauscht ($x^5$ müsste eingeschränkt werden): alle drei Zeilen 0.}
\end{lpaufgabe}

\begin{lpaufgabe}{G7}{$k(x) = 2\sqrt{x+3} - 4$ · $m(x) = \sqrt[3]{x+1} - 2$}{4 P}
\begin{raster}
(a) Definitionsmenge und Startpunkt & 1\A & Radikand $\geq 0$: $x + 3 \geq 0 \Rightarrow D = [-3;\, +\infty[$ (gleichwertig $x \geq -3$); $S(-3 \mid -4)$. Beides nötig.\\
(a) Nullstelle & 1\E & $2\sqrt{x+3} = 4 \Rightarrow \sqrt{x+3} = 2 \Rightarrow x + 3 = 4 \Rightarrow x_0 = 1$\\
(b) Definitionsmenge & 1\A & $D = \mathbb{R}$\\
(b) Begründung & 1 & Der Wurzelexponent $3$ ist ungerade, darum darf der Radikand negativ sein: $\sqrt[3]{-8} = -2$. Bei $k$ ist er gerade, dort muss der Radikand $\geq 0$ bleiben.\\
\end{raster}
\fehler{$D = [3;\, +\infty[$ (Vorzeichen von $u$): Zeile (a) erste 0, auch wenn der Startpunkt $(3 \mid -4)$ dazu passt $\to$ 3 von 4\,P.
$S(-3 \mid 4)$ bei richtiger Definitionsmenge: Zeile (a) erste 0 $\to$ 3 von 4\,P.
$x_0 = 13$ (die $2$ vor der Wurzel nicht weggeteilt, also $\sqrt{x+3} = 4$ gerechnet): Nullstellen-Zeile 0 $\to$ 3 von 4\,P.
$x_0 = 5$ (quadriert, ohne die $2$ mitzuquadrieren: $2(x+3) = 16$): Nullstellen-Zeile 0 $\to$ 3 von 4\,P.
$x_0 = -7$ (Vorzeichen beim Auflösen): Nullstellen-Zeile 0 $\to$ 3 von 4\,P. Die Probe $k(-7)$ ist in $D$ gar nicht möglich — ein Hinweis darauf ist richtig, zählt aber nicht zusätzlich.
Bei (b) $D = [-1;\, +\infty[$ (wie bei einer geraden Wurzel): Zeile (b) erste 0; eine Begründung, die dazu passt, ist ebenfalls falsch $\to$ höchstens 2 von 4\,P.
\emph{Hinweis für die Bewertung:} Wer bei (b) $D = \mathbb{R}_0^+$ schreibt und sich dabei auf die Konvention der Themenseite 3.2b beruft, hat die Abweichung verstanden — dann beide (b)-Zeilen voll.}
\end{lpaufgabe}

\begin{lpaufgabe}{G8}{Der Grösse nach}{2 P}
\begin{raster}
Reihenfolge & 1\A & $\sqrt{x} > x > x^{2}$ für $0 < x < 1$; zählt auch ohne Rechenweg\\
Zahlenbeispiel & 1 & z.\,B. $x = 0.25$: $\sqrt{0.25} = 0.5$, $x = 0.25$, $x^{2} = 0.0625$ — alle drei Werte nötig\\
\end{raster}
\fehler{$x^{2} > x > \sqrt{x}$ (Fall $x > 1$ genommen): Reihenfolge 0; ein Beispiel aus $0 < x < 1$ kann diese Reihenfolge nicht belegen, und ein Beispiel mit $x > 1$ beantwortet die Frage nicht $\to$ 0 von 2\,P.
Reihenfolge richtig, Zahlenbeispiel fehlt oder nennt nur einen Wert: zweite Zeile 0 $\to$ 1 von 2\,P.
Als Beispiel ein $x > 1$ gewählt: zweite Zeile 0, weil es die Aussage nicht belegt $\to$ 1 von 2\,P.}
\end{lpaufgabe}

\needspace{14\baselineskip}
\section*{4 · Selbsteinschätzung}
\begin{tabularx}{\linewidth}{@{}l X@{}}\toprule
23 – 26 P & Die geprüften Teile sitzen. Wo du Punkte verloren hast: das Kapitel dieser Aufgabe nochmals (Zuordnung unten).\\
18 – 22 P & Den schwächsten Teil nochmals: Simulation und Übungen des Kapitels, in dem du die meisten Punkte verloren hast.\\
12 – 17 P & Zurück zu den Kapiteln aller Aufgaben, in denen du Punkte verloren hast.\\
0 – 11 P & Zurück zu Kapitel 1 und von dort der Reihe nach weiter.\\\midrule
\multicolumn{2}{@{}p{\linewidth}@{}}{\small\textbf{Aufgabe $\to$ Kapitel:} G1, G2 $\to$ Kapitel 1; G3 $\to$ Kapitel 2; G4 $\to$ Kapitel 3 (Nullstellen und Asymptoten); G5, G6 $\to$ Kapitel 4; G7 $\to$ Kapitel 5 (beide Paritäten); G8 $\to$ Kapitel 5}\\\bottomrule
\end{tabularx}
\end{document}
